% !TEX root = gnm.tex
% !TEX recipe = LuaLaTeX+se

\chapter[Lent]{TEMPUS QUADRAGESIMAE}

The traditional Gospel acclaimation for Lent is the \emph{Tract}.
The Missal of Paul VI specifies that the Gospel acclaimation from
the lectionary is to be used in Lent, but also says that it is also
``possible to sing another Psalm or Tract, as found in the Graduale.''
It is not clear if this allows for the use of a Tract that is \emph{not}
found (but is derived from) the Graduale---e.g., a translation or
abridged version (as provided here.) Regardless, it is not an
unheard of practice to substitute a variant of the Tract for the
Gospel acclaimation in the OF, and as of this writing, the author
knows of no correction by the bishops.

In the spirit of interpreting
ambiguities in the law in favor of freedom---which, indeed, seems to
have been the intent of Vatican II, along with the restoration
of Gregorian Chant in general---this author is of the opinion that
the Tracts given in this book may be used. However, we have also 
provided settings for the acclaimations from the lectionary as well,
based on other graduals in the repertoire.
Which---or even both---to use should be discussed with one's pastor.


\section{First Sunday of Lent}

\gregorioscore{Acclaimations/one-does-not/drj-mode-ii}


